\documentclass[12pt,a4paper]{report}
\usepackage[]{amsmath,amsthm,amssymb,amscd}
\usepackage[latin1]{inputenc}
%\usepackage{fullpage}

\topmargin 0pt
\advance \topmargin by -\headheight
\advance \topmargin by -\headsep
     
\textheight 246mm
     
\oddsidemargin 0pt
\evensidemargin \oddsidemargin
\marginparwidth 0.5in
     
\textwidth 159mm


%               Theorem Environments
\theoremstyle{definition}
\newtheorem{ex}{}
\newcommand{\pd}[1]{\frac{\partial}{\partial #1}} %\pd{x}
%               Subquestions numbered with letters
\renewcommand{\theenumi}{\textbf{\alph{enumi}}}

%               Maths definitions

\newcommand{\NN}{\ensuremath{\mathbb N}}
\newcommand{\ZZ}{\ensuremath{\mathbb Z}}
\newcommand{\CC}{\ensuremath{\mathbb C}}
\newcommand{\RR}{\ensuremath{\mathbb R}}
\newcommand{\TT}{\ensuremath{\mathbb T}}
\newcommand{\g}{\ensuremath{\frak{g}}}
\newcommand{\h}{\ensuremath{\frak{h}}}
\newcommand{\ft}{\ensuremath{\frak{t}}}
\newcommand{\cB}{\mathcal{B}}
\newcommand{\cN}{\mathcal{N}}
\newcommand{\cL}{\mathcal{L}}
\newcommand{\cD}{\mathcal{D}}
%\newcommand{\pd}[1]{\frac{\partial}{\partial #1}} %\pd{x}
%\newcommand{\pd}[1]{\frac{\partial}{\partial #1}} %\pd{x}

\begin{document}
\noindent
{\sffamily\large Marco Zambon 
\hfill   a discutir:   Jueves, 3/05/2012\\Geometr\'ia III, curso 2011-12}
 

\noindent
\hrulefill{}\\
\begin{center}\bfseries\large\textsf{Hoja 5    \vspace{-2mm}}
\end{center}

\noindent
%\hrulefill{}\\
 
 

\begin{ex}
\begin{enumerate}
\item   Sea $A\in Mat(n,\RR)$ una matriz cuadrada. Para cada $t\in \RR$ obtenemos un difeomorfismo 
$e^{tA} \colon \RR^n \to \RR^n$, dado por $v \mapsto e^{tA} v$ donde $e^{tA}:=\sum_{i=0}^{\infty}\frac{(tA)^i}{i!}$ es la exponencial de matrices.
  Encuentra un campo vectorial en $\RR^n$ cuyo flujo $F_t$ sea dado por $e^{tA}$.
  
\item  Tomamos $A=\left(\begin{array}{ccc}0 & 1  \\-1 & 0\end{array}\right) \in Mat(2,\RR)$. 
Calcula explicitamente $e^{tA}$ y el campo vectorial en $\RR^2$ cuyo flujo es dado por $e^{tA}$.
\end{enumerate}
\end{ex}

 \begin{ex}
  Considera $\RR^2$ con coordenadas $x,y$. Calcula los siguientes crochetes de Lie  de campos vectoriales:
\begin{enumerate}
\item $[x\pd{y},x\pd{y}]$ 
\item $[x\pd{y},\pd{x}]$
\end{enumerate}
\end{ex}


\begin{ex} 
Sea $n\ge 1$. Demuestra que $Mat(n,\RR)$, con $[A,B]:=AB-BA$, es un \'algebra de Lie. 
\end{ex}


 \begin{ex}
  Sea $V$ un espacio vectorial y $\{e_1,\dots,e_n\}\subset V$ una base de $V$. Sea $\{f_1,\dots,f_n\}\subset V^*$ la base dual. Dados vectores
$v=\sum_{i=1}^n a_i e_i$ y $w=\sum_{i=1}^n b_i e_i$ de $V$ (aqu\'i $a_i,b_i\in \RR$), calcula
$$(f_1\wedge f_2)(v,w).$$
 
\end{ex}

 \begin{ex} Sea $V$ un espacio vectorial y $\{f_1,\dots,f_n\}\subset V^*$ una base de $V^*$. Demuestra,  para cada $k\in\{1,\dots,n\}$, que
 $$\{f_{i_1}\wedge\dots\wedge f_{i_k}: 1\le i_1 <\dots<i_k\le n\}$$
 es una base de $\wedge^k V^*$.
\end{ex}

 \begin{ex}
 Considera $\RR^3$ con coordenadas $x,y,z$. 
 \begin{enumerate}
 
\item Calcula $(dx \wedge dy \wedge dz) (\pd{y},\pd{x}, x \pd{z})$
\item Encuentra 2 campos vectoriales $X$ y $Y$ en $\RR^3$, linealmente independientes en cada punto, tal que
$$(xdy+dz)(X)=(xdy+dz)(Y)=0.$$
\end{enumerate}
\end{ex}
 \end{document}
 %%%%%%% 
 
 

 \begin{ex}
\begin{enumerate}
\item 
\end{enumerate}
\end{ex}
 
 
 
 